% ============================================================================
% Appendix: Phase 2 NMCP — Non-Monotone Core-Profile Provenance Test
% ============================================================================
% Status: Provenance-first test of the NMCP loophole
% Scope: Test whether any admissible physical mechanism within S_EH + S_NG
%        breaks the q=0-centered monotone sign logic established in N7
% Phase 2 / NMCP WP2 deliverable
% Governing documents: audit/NMCP_WP1_DONOR_NORMALIZATION.md,
%                      audit/MINIMAL_CLASS_CLOSURE_MEMO_AFTER_N1_N7_N2.md
% ============================================================================

\chapter{Non-Monotone Core-Profile Provenance Test for the Node-Well Sector}
\label{app:P2_NMCP}

% ============================================================================
\section{Scope and Framing}
\label{sec:NMCP:scope}
% ============================================================================

This appendix is a \textbf{provenance-first test} of the non-monotone
core-profile (NMCP) loophole identified in Appendix~\ref{app:P2_N7_core}
(\S Escape Route). It does not implement a non-monotone profile and
scan for metastability. It asks the prior question: \emph{can any
admissible physical mechanism within $S_{\mathrm{EH}} + S_{\mathrm{NG}}$
generate an effective core energy that is not minimized at
$q = 0$ (Steiner)?}

\textbf{Why provenance first:} The NMCP lane has three freely tunable
parameters (peak position $q_*$, width $w$, amplitude $A$) that can
produce any desired $\VB$. Without provenance---a physical mechanism
that selects the non-monotone structure before the target is
known---any numerical success is indistinguishable from a
phenomenological node well relabeled as ``core physics.''

\textbf{What this appendix tests:}
\begin{itemize}
    \item Whether any of the three WP1 candidate mechanisms
          (\S\ref{sec:NMCP:candidates}) breaks the N7 sign argument
    \item Whether the sign argument extends beyond the monotone/separable
          class
    \item Whether non-monotonicity has provenance within
          $S_{\mathrm{EH}} + S_{\mathrm{NG}}$ or requires
          non-minimal extensions
\end{itemize}

\textbf{What this appendix does not claim:}
\begin{itemize}
    \item It does not compute a specific non-monotone $f(q/\delta)$.
    \item It does not scan profiles or fit parameters.
    \item It does not derive $V_B$, $\tau_n$, or any downstream quantity.
    \item A no-go is an allowed and expected outcome.
\end{itemize}

% ============================================================================
\section{Governing Constraint from N7}
\label{sec:NMCP:constraint}
% ============================================================================

\subsection{The Monotone Profile No-Go}

Appendix~\ref{app:P2_N7_core}, Theorem~\ref{thm:monotone_nogo}
established:

\begin{resultbox}[title=N7 Monotone Profile No-Go (restated)]
If the core profile $f(q/\delta)$ satisfies
(i)~$f(0) = \max f$ and (ii)~$f'(u) \leq 0$ for $u > 0$,
then $V(q) = V_{\mathrm{geom}}(q) + V_{\mathrm{core}}(q)$
has no secondary minimum: $V'(q) > 0$ for all $q > 0$.
\tagDc
\end{resultbox}

This eliminated the natural model class (all standard profiles:
Gaussian, Lorentzian, exponential, disk).

\subsection{The Sign Argument}

Underlying the monotone no-go is a \textbf{physical sign argument}
(Appendix~\ref{app:P2_N7_core}, \S\ref{sec:P2N7:reduction}):

\begin{enumerate}
    \item \textbf{Binding energy} (negative contribution): the three
          worldsheets merge in the core, saving energy relative to
          separate arms. At $q = 0$ (Steiner, $\Zthree$-symmetric),
          the overlap is \emph{maximal}. Any displacement $q > 0$
          reduces the symmetry ($\Zthree \to \Ztwo$) and reduces the
          overlap, making binding less negative:
          \begin{equation}
              E_{\mathrm{bind}}(0) \leq E_{\mathrm{bind}}(q)
              \quad \text{for } q > 0
              \label{eq:NMCP_bind}
          \end{equation}

    \item \textbf{Strain energy} (positive contribution): stress
          concentration at the vertex costs energy. For an elastic
          medium, the strain energy density scales as $\sigma^2 / Y$.
          By convexity of $x \mapsto x^2$, equal stress distribution
          ($\Zthree$-symmetric at $q = 0$) minimizes total strain
          energy:
          \begin{equation}
              E_{\mathrm{strain}}(0) \leq E_{\mathrm{strain}}(q)
              \quad \text{for } q > 0
              \label{eq:NMCP_strain}
          \end{equation}
\end{enumerate}

Both contributions favor $q = 0$. Combined:
\begin{equation}
    V_{\mathrm{core}}(0) = E_{\mathrm{bind}}(0) + E_{\mathrm{strain}}(0)
    \leq E_{\mathrm{bind}}(q) + E_{\mathrm{strain}}(q)
    = V_{\mathrm{core}}(q)
    \label{eq:NMCP_sign_combined}
\end{equation}

\textbf{Epistemic status:} \tagDcI. Uses (i)~separable ansatz [P],
(ii)~elastic-medium analogy [P], (iii)~convexity of $\sigma^2$ [Der].

\textbf{The question for NMCP:} Does any admissible mechanism
invalidate Eq.~\eqref{eq:NMCP_sign_combined}?

% ============================================================================
\section{Candidate Admissible Mechanisms}
\label{sec:NMCP:candidates}
% ============================================================================

The NMCP WP1 normalization (\texttt{audit/NMCP\_WP1\_DONOR\_NORMALIZATION.md},
\S7.1) identified exactly three candidate mechanisms that could, in
principle, break the monotone sign logic. Each must be tested against:
\begin{itemize}
    \item[(A)] the sign argument (which assumption does it break?)
    \item[(B)] the admissibility criteria (WP1 \S7.1, items 1--5)
    \item[(C)] the quick-falsification conditions (WP1 \S11, QF-1 through QF-7)
\end{itemize}

\subsection{Candidate 1: $\Zthree \to \Ztwo$ Internal Symmetry Transition}

\textbf{Mechanism:} At some critical displacement $q_*$, the
three-fold symmetric junction core reorganizes to a two-fold
configuration, releasing stress energy.

\textbf{How it would break the sign argument:}
The strain energy inequality Eq.~\eqref{eq:NMCP_strain} relies on
$\Zthree$ being the minimum-strain configuration for all $q$. If
a $\Zthree \to \Ztwo$ transition occurs at $q_*$, the strain energy
could jump \emph{downward} at $q_*$ because the two-fold configuration
might distribute stress more efficiently for asymmetric geometries.

\textbf{Provenance test:}
\begin{enumerate}
    \item \textbf{What drives the transition?} Within the Nambu--Goto
          action, the worldsheet energy depends only on total area (or
          total length in the 1D reduction). There is no order parameter
          that distinguishes $\Zthree$ from $\Ztwo$ core structure.
          The Nambu--Goto action is proportional to arm lengths---it has
          no internal degrees of freedom at the vertex.

    \item \textbf{What about the regularized core?} The regularization
          replaces the distributional vertex with a smooth merging region
          of width $\delta$. Within an elastic-medium model, the stress
          distribution evolves \emph{continuously} with $q$. There is no
          bifurcation mechanism: as $q$ increases, the stress on arm 1
          (toward $\mathbf{P}_1$) decreases while stress on arms 2--3
          increases. The $\Ztwo$ residual symmetry ($\mathbf{P}_2
          \leftrightarrow \mathbf{P}_3$) is maintained for all $q$, but
          this is a continuous reduction from $\Zthree$, not a
          first-order transition.

    \item \textbf{First-order transition requirement:} A genuine
          $\Zthree \to \Ztwo$ transition would require either:
          \begin{itemize}
              \item[(a)] a first-order phase transition in the core
                    material (requiring a free energy with two competing
                    local minima as a function of internal order parameter),
                    or
              \item[(b)] an instability threshold (a mode going soft at
                    critical $q$, leading to spontaneous symmetry breaking
                    within the core).
          \end{itemize}
          Neither exists in $S_{\mathrm{EH}} + S_{\mathrm{NG}}$. The
          Nambu--Goto action has no internal degrees of freedom that could
          undergo a phase transition. The Einstein--Hilbert action couples
          to the brane stress-energy but does not introduce order parameters
          at the junction vertex.
\end{enumerate}

\textbf{Verdict on Candidate 1:}
\begin{resultbox}[title=Candidate 1: $\Zthree \to \Ztwo$ Transition]
\textbf{Classification:} No provenance within
$S_{\mathrm{EH}} + S_{\mathrm{NG}}$.

The $\Zthree \to \Ztwo$ transition requires internal core dynamics
(phase transition or instability) that are not present in the minimal
action class. The Nambu--Goto action has no internal vertex degrees of
freedom; the stress distribution evolves continuously with $q$; no
bifurcation mechanism exists.

This mechanism would require non-minimal extensions (e.g., additional
fields coupling to the junction vertex, or higher-order brane terms
with spontaneous symmetry-breaking potential).

\textbf{Quick-falsification:} QF-6 triggered --- the mechanism requires
physics not present in $S_{\mathrm{EH}} + S_{\mathrm{NG}}$.
\tagDc
\end{resultbox}

\subsection{Candidate 2: Non-Separable Transverse-Longitudinal Coupling}

\textbf{Mechanism:} If $\rho_{\mathrm{core}}(r_\perp, q) \neq
g_\perp(r_\perp/r_0) \times f(q/\delta)$, a cross-coupling between
transverse and longitudinal degrees of freedom could make the
effective $f_{\mathrm{eff}}(q)$ (obtained after transverse integration)
non-monotone.

\textbf{How it would break the sign argument:}
The separability assumption [P] was one of the stated premises of the
sign argument. If separability fails, the binding and strain energy
contributions cannot be factored into independent transverse and
longitudinal parts. The effective $q$-dependence after integrating
out $r_\perp$ could, in principle, differ from any monotone function.

\textbf{Provenance test:}
\begin{enumerate}
    \item \textbf{Does the sign argument depend on separability?}
          Consider the total core energy as a function of $q$, without
          assuming separability:
          \begin{equation}
              V_{\mathrm{core}}(q) = \int d^2 r_\perp\;
              \mathcal{E}_{\mathrm{core}}(r_\perp, q)
              \label{eq:NMCP_nonsep_Vcore}
          \end{equation}
          where $\mathcal{E}_{\mathrm{core}}$ is the local core energy
          density.

    \item \textbf{Binding energy at each $r_\perp$:}
          At any fixed transverse position $r_\perp$ within the core,
          the binding energy density depends on how efficiently the three
          worldsheets merge at that location. At $q = 0$ (Steiner), the
          merger is $\Zthree$-symmetric at every $r_\perp$. At $q > 0$,
          the asymmetry propagates into the core interior: the sheets
          are closer to arm-1 at every $r_\perp$, reducing the symmetric
          merger efficiency.

          Therefore, at each $r_\perp$:
          \begin{equation}
              \mathcal{E}_{\mathrm{bind}}(r_\perp, 0)
              \leq \mathcal{E}_{\mathrm{bind}}(r_\perp, q)
              \quad \text{for } q > 0
              \label{eq:NMCP_bind_local}
          \end{equation}

    \item \textbf{Strain energy at each $r_\perp$:}
          By the same convexity argument ($\sigma^2$ is convex), the
          strain energy density at each $r_\perp$ is minimized by the
          symmetric stress distribution at $q = 0$:
          \begin{equation}
              \mathcal{E}_{\mathrm{strain}}(r_\perp, 0)
              \leq \mathcal{E}_{\mathrm{strain}}(r_\perp, q)
              \quad \text{for } q > 0
              \label{eq:NMCP_strain_local}
          \end{equation}

    \item \textbf{Integration preserves the inequality:}
          Integrating Eqs.~\eqref{eq:NMCP_bind_local} and
          \eqref{eq:NMCP_strain_local} over all $r_\perp$:
          \begin{equation}
              V_{\mathrm{core}}(0) = \int d^2 r_\perp\;
              [\mathcal{E}_{\mathrm{bind}} + \mathcal{E}_{\mathrm{strain}}]
              (r_\perp, 0)
              \leq V_{\mathrm{core}}(q)
              \label{eq:NMCP_nonsep_sign}
          \end{equation}
          This holds \textbf{without the separability assumption}. The
          sign argument is \emph{local in $r_\perp$}---it applies at each
          transverse position independently. Separability was used in N7
          for mathematical convenience (to extract the factored form
          $V_{\mathrm{core}} = -E_0 f$), but the physical conclusion
          ($q = 0$ is the core energy minimum) does not depend on it.
\end{enumerate}

\textbf{Verdict on Candidate 2:}
\begin{resultbox}[title=Candidate 2: Non-Separable Coupling]
\textbf{Classification:} Sign argument extends to the non-separable
case. No escape route.

The binding and strain energy arguments that establish $q = 0$ as the
core energy minimum are \emph{local} in the transverse coordinate
$r_\perp$. They apply at each $r_\perp$ independently, and integration
over $r_\perp$ preserves the inequality. The non-separable case
inherits the same sign logic as the separable case.

Non-separable transverse-longitudinal coupling can modify the
\emph{magnitude} and \emph{shape} of $V_{\mathrm{core}}(q)$ relative
to the separable approximation, but it cannot change the sign:
$V_{\mathrm{core}}(0) \leq V_{\mathrm{core}}(q)$ for all $q > 0$.

\textbf{Quick-falsification:} QF-5 triggered --- the sign logic holds
for all physically reasonable core models, including non-separable ones.
\tagDc
\end{resultbox}

\subsection{Candidate 3: Core Topology Change}

\textbf{Mechanism:} At some critical displacement $q_*$, the core
topology changes (e.g., from connected to disconnected junction),
releasing topological energy.

\textbf{How it would break the sign argument:}
The sign argument assumes continuous deformation. A topological
transition could produce a non-analytic (discontinuous) change in
$V_{\mathrm{core}}$ at $q_*$, bypassing the continuous sign logic.

\textbf{Provenance test:}
\begin{enumerate}
    \item \textbf{What constitutes a topology change for a Y-junction?}
          The junction is a topological defect where three worldsheets
          meet. A topology change would mean the defect splits into
          components, or two arms merge into one. In either case, this
          is a singular event---the junction either exists or it does not.

    \item \textbf{Is there a framework for this within
          $S_{\mathrm{EH}} + S_{\mathrm{NG}}$?}
          The Nambu--Goto action describes worldsheets with fixed
          topology. Topology changes (worldsheet splitting, merging)
          are not part of the classical action---they require either
          quantum path integral contributions (worldsheet instantons)
          or additional topological terms not present in $S_{\mathrm{NG}}$.

    \item \textbf{Physical plausibility:} A Y-junction displaced by
          $q \ll R$ does not approach any topological transition. The
          relevant displacement range ($q \sim \delta \approx 0.1\fm$,
          much smaller than $R \sim 1\fm$) is far from any singular
          configuration.
\end{enumerate}

\textbf{Verdict on Candidate 3:}
\begin{resultbox}[title=Candidate 3: Core Topology Change]
\textbf{Classification:} No provenance within
$S_{\mathrm{EH}} + S_{\mathrm{NG}}$. Highly speculative.

The Nambu--Goto action has fixed worldsheet topology. Topology changes
require quantum effects or additional topological terms not present
in the minimal action class. The relevant displacement range
($q \sim \delta \ll R$) is far from any singular configuration.

No framework exists for computing topology-change contributions to
$V(q)$ within $S_{\mathrm{EH}} + S_{\mathrm{NG}}$.

\textbf{Quick-falsification:} QF-6 triggered --- requires physics not
present in $S_{\mathrm{EH}} + S_{\mathrm{NG}}$.
\tagDc
\end{resultbox}

% ============================================================================
\section{Derivation: Extension of the Sign Argument Beyond Separability}
\label{sec:NMCP:derivation}
% ============================================================================

The key new result of this appendix is the extension of the N7 sign
argument to the \textbf{general (non-separable) case}. This closes the
non-separable loophole that was left open by the N7 monotone no-go
theorem.

\begin{theorem}[General Core Sign Theorem]
\label{thm:NMCP_general_sign}
Let the junction core energy be
\begin{equation}
    V_{\mathrm{core}}(q) = \int d^2 r_\perp\;
    \mathcal{E}_{\mathrm{core}}(r_\perp, q)
    \label{eq:NMCP_general_Vcore}
\end{equation}
where $\mathcal{E}_{\mathrm{core}}(r_\perp, q)$ is the local core
energy density at transverse position $r_\perp$ and node displacement
$q$. Assume:
\begin{enumerate}
    \item[(A)] \textbf{Elastic brane:} The brane stress-energy in the
          core region is described by an elastic medium with tension
          $\sigma$ and stress distribution $\{\sigma_i(r_\perp, q)\}_{i=1}^3$
          along the three arm directions.
    \item[(B)] \textbf{$\Zthree$-symmetric equilibrium:} At $q = 0$
          (Steiner), the stress distribution is $\Zthree$-symmetric:
          $\sigma_1(r_\perp, 0) = \sigma_2(r_\perp, 0)
          = \sigma_3(r_\perp, 0)$ for all $r_\perp$.
    \item[(C)] \textbf{Convex strain energy:} The strain energy density
          $\varepsilon_{\mathrm{strain}} = \sum_i \sigma_i^2 / (2Y)$
          is convex in the stress components $\{\sigma_i\}$ at fixed
          total stress $\sum_i \sigma_i = \mathrm{const}$.
    \item[(D)] \textbf{Overlap monotonicity:} The binding energy density
          at each $r_\perp$ satisfies
          $\mathcal{E}_{\mathrm{bind}}(r_\perp, 0) \leq
          \mathcal{E}_{\mathrm{bind}}(r_\perp, q)$ for all $q > 0$.
          (The $\Zthree$-symmetric configuration maximizes the binding
          efficiency of the three-sheet merger.)
\end{enumerate}
Then:
\begin{equation}
    V_{\mathrm{core}}(0) \leq V_{\mathrm{core}}(q)
    \quad \text{for all } q > 0
    \label{eq:NMCP_general_sign_result}
\end{equation}
The Steiner configuration minimizes the core energy.
\end{theorem}

\begin{proof}
At each $r_\perp$, the core energy density is:
\begin{equation}
    \mathcal{E}_{\mathrm{core}}(r_\perp, q)
    = \mathcal{E}_{\mathrm{bind}}(r_\perp, q)
    + \mathcal{E}_{\mathrm{strain}}(r_\perp, q)
\end{equation}

By assumption (D):
$\mathcal{E}_{\mathrm{bind}}(r_\perp, 0) \leq
\mathcal{E}_{\mathrm{bind}}(r_\perp, q)$.

By assumptions (B) and (C): the strain energy density at $q = 0$ is
minimized by the $\Zthree$-symmetric stress distribution
(equal distribution minimizes $\sum \sigma_i^2$ at fixed
$\sum \sigma_i$, by convexity). For $q > 0$, the stress distribution
is asymmetric, so
$\mathcal{E}_{\mathrm{strain}}(r_\perp, 0) \leq
\mathcal{E}_{\mathrm{strain}}(r_\perp, q)$.

Adding the two inequalities at each $r_\perp$:
\begin{equation}
    \mathcal{E}_{\mathrm{core}}(r_\perp, 0)
    \leq \mathcal{E}_{\mathrm{core}}(r_\perp, q)
    \quad \text{for all } r_\perp,\; q > 0
\end{equation}

Integrating over $r_\perp$:
\begin{equation}
    V_{\mathrm{core}}(0) = \int d^2 r_\perp\;
    \mathcal{E}_{\mathrm{core}}(r_\perp, 0)
    \leq \int d^2 r_\perp\;
    \mathcal{E}_{\mathrm{core}}(r_\perp, q)
    = V_{\mathrm{core}}(q)
\end{equation}
\end{proof}

\textbf{What this theorem adds beyond N7:}
\begin{itemize}
    \item N7 Theorem~\ref{thm:monotone_nogo} required the separable
          ansatz and the monotone profile condition.
    \item Theorem~\ref{thm:NMCP_general_sign} works for
          \textbf{any} core energy density, separable or not.
    \item The crucial insight: the sign argument is \emph{local in
          $r_\perp$}. Separability was a mathematical convenience,
          not a physical requirement.
\end{itemize}

\textbf{Assumptions:}
The theorem rests on assumptions (A)--(D), all tagged:
\begin{itemize}
    \item[(A)] Elastic brane: [P]. This is the same model-class
          assumption as N7. The brane is treated as an elastic medium.
    \item[(B)] $\Zthree$-symmetric equilibrium: [Dc]. Follows from the
          equilateral boundary configuration
          (Appendix~\ref{app:vq_chosen_path}, boundary points
          Eq.~\eqref{eq:boundary_points}).
    \item[(C)] Convex strain energy: [Der]. $x \mapsto x^2$ is convex;
          equal distribution minimizes $\sum x_i^2$ at fixed $\sum x_i$
          (Schur convexity / rearrangement inequality).
    \item[(D)] Overlap monotonicity: [P]. Physically motivated: symmetric
          merger of three sheets at the vertex produces maximal binding
          efficiency. But this is a model assumption, not a theorem.
\end{itemize}

\textbf{The weakest assumption is (D)} --- overlap monotonicity.
Breaking this assumption would require a mechanism in which the
binding of three worldsheets becomes \emph{more efficient} when the
node is displaced from the symmetric position. This would contradict
the generic principle that symmetric configurations maximize overlap
for identical components. No such mechanism has been identified within
$S_{\mathrm{EH}} + S_{\mathrm{NG}}$.

% ============================================================================
\section{Provenance Test}
\label{sec:NMCP:provenance}
% ============================================================================

The WP1 provenance test (\S7.3) requires:

\begin{quote}
\emph{Remove all knowledge of $V_B$, $\tau_n$, $\Delta m_{np}$, and
the Variant~3 Gaussian. Does the non-monotone profile still emerge
from the stated physics?}
\end{quote}

\textbf{Answer: No.}

\begin{enumerate}
    \item \textbf{Is the off-center extremum generated by any tested
          mechanism?}
          No. All three WP1 candidate mechanisms fail:
          \begin{itemize}
              \item Candidate 1 ($\Zthree \to \Ztwo$ transition):
                    requires non-minimal physics (QF-6)
              \item Candidate 2 (non-separable coupling):
                    sign argument extends to general case (QF-5)
              \item Candidate 3 (topology change):
                    requires non-minimal physics (QF-6)
          \end{itemize}

    \item \textbf{Is non-monotonicity effectively inserted by choice of
          profile family?}
          The only way to obtain $f(q_*/\delta) > f(0)$ within the
          current framework is to postulate such a profile. No physical
          mechanism generates it. Therefore, any non-monotone profile
          used in a WP2 computation would be an inserted function,
          not a derived one.

    \item \textbf{Does the mechanism produce non-monotonicity before
          fitting?}
          No mechanism produces non-monotonicity at all. The sign
          argument (now extended to non-separable case,
          Theorem~\ref{thm:NMCP_general_sign}) establishes that
          $V_{\mathrm{core}}(0) \leq V_{\mathrm{core}}(q)$ under the
          elastic brane model assumptions. Non-monotonicity requires
          breaking these assumptions.

    \item \textbf{Does any new parameter simply play the role of hidden
          $q_*$ / width / amplitude?}
          Not applicable---no mechanism survived the provenance test.
          If one were forced to proceed (against the evidence), any
          non-monotone profile would introduce exactly the three
          freely tunable parameters ($q_*$, $w$, amplitude) flagged
          in WP1 \S9 as the primary circularity risks.
\end{enumerate}

\textbf{Provenance test result: FAILED.} No admissible mechanism within
$S_{\mathrm{EH}} + S_{\mathrm{NG}}$ generates non-monotone core
energy. The non-monotone loophole is formally open (the general sign
theorem rests on assumptions (A)--(D), which are model-dependent) but
physically empty within the current framework.

% ============================================================================
\section{Outcome Classification}
\label{sec:NMCP:outcome}
% ============================================================================

\begin{resultbox}[title=NMCP WP2 Outcome: Bounded No-Go]
\textbf{Classification:} Bounded no-go within
$S_{\mathrm{EH}} + S_{\mathrm{NG}}$ under the elastic brane model.

\textbf{What was established:}
\begin{enumerate}
    \item \textbf{General sign theorem [Dc]:} The core energy
          $V_{\mathrm{core}}(q)$ is minimized at $q = 0$ (Steiner)
          for \emph{any} core energy density (separable or not),
          provided the elastic brane model assumptions (A)--(D) hold.
          This extends the N7 monotone no-go from the separable class
          to the general case.

    \item \textbf{All three WP1 candidate mechanisms rejected:}
          \begin{itemize}
              \item $\Zthree \to \Ztwo$ transition: no provenance
                    within $S_{\mathrm{EH}} + S_{\mathrm{NG}}$ (QF-6)
              \item Non-separable coupling: sign argument extends to
                    general case (QF-5)
              \item Core topology change: no provenance within
                    $S_{\mathrm{EH}} + S_{\mathrm{NG}}$ (QF-6)
          \end{itemize}

    \item \textbf{Provenance test failed:} No physical mechanism
          generates non-monotone core energy within the minimal action
          class. Non-monotonicity can only be inserted by postulating
          a profile --- which is a disguised phenomenological node well.
\end{enumerate}

\textbf{What this means:}
The NMCP loophole is \textbf{physically empty} within
$S_{\mathrm{EH}} + S_{\mathrm{NG}}$. The non-monotone lane survives
only as a formal mathematical possibility (the general sign theorem
has model-dependent assumptions), but no mechanism within the minimal
action class can instantiate it.

\textbf{Scope of the no-go:}
\begin{itemize}
    \item Any core energy density (separable or non-separable)
    \item Any regularization model satisfying the elastic brane
          assumptions (A)--(D)
    \item Any values of $\sigma$, $\delta$, $L_0$, $R$
\end{itemize}

\textbf{What could break the no-go:}
\begin{itemize}
    \item A mechanism that violates assumption (D) --- overlap
          monotonicity. This would require asymmetric binding to be
          \emph{more efficient} than symmetric binding, which contradicts
          the generic principle for identical components.
    \item Non-minimal action extensions: additional fields, higher-order
          brane terms, topological contributions. These are outside
          $S_{\mathrm{EH}} + S_{\mathrm{NG}}$.
\end{itemize}

\textbf{Why ``bounded no-go'' not ``absolute no-go'':}
The theorem rests on the elastic brane model (A)--(D). A fundamentally
different model of the junction core (not elastic, not stress-based)
could evade the theorem. But no such model exists within
$S_{\mathrm{EH}} + S_{\mathrm{NG}}$, and the elastic-medium treatment
is the natural reduction of the Nambu--Goto brane.

\textbf{Epistemic tag:} \tagDc\ (structural no-go within model class).
\end{resultbox}

% ============================================================================
\section{What the Route Can and Cannot Claim}
\label{sec:NMCP:claims}
% ============================================================================

\textbf{The NMCP route may claim:}
\begin{itemize}
    \item That it tested all three WP1-identified candidate mechanisms
    \item That it extended the N7 sign argument beyond separability
    \item That the provenance test was applied honestly
    \item That the result is a genuine bounded no-go, not a failure
          to compute
\end{itemize}

\textbf{The NMCP route may not claim:}
\begin{itemize}
    \item That non-monotone core profiles are absolutely impossible
          (the theorem has model-dependent assumptions)
    \item That the double-well is falsified (only the NMCP loophole
          within $S_{\mathrm{EH}} + S_{\mathrm{NG}}$ is closed)
    \item Any back-promotion of neutron-line epistemic status
    \item That the remaining non-minimal escape routes (higher-order
          brane terms, topological contributions) are also closed
\end{itemize}

% ============================================================================
\section{Numerical Verification}
\label{sec:NMCP:numerical}
% ============================================================================

The verification code (\texttt{code/p2\_nmcp\_provenance\_check.py})
confirms the analytic results:

\begin{enumerate}
    \item \textbf{Sign check:} For a comprehensive family of core
          energy models (Gaussian, Lorentzian, exponential, uniform,
          and their non-separable generalizations), $V_{\mathrm{core}}(0)
          \leq V_{\mathrm{core}}(q)$ for all sampled $q > 0$.
    \item \textbf{Combined potential check:} $V(q) =
          V_{\mathrm{geom}}(q) + V_{\mathrm{core}}(q)$ has
          $V'(q) > 0$ for all $q > 0$ for all tested models.
    \item \textbf{Non-monotone forced test:} If a non-monotone
          profile is manually inserted (e.g., $f(u) = 1 + A\,u^2
          e^{-u^2}$ with $A > 0$, which has a maximum at $u > 0$),
          the code correctly detects metastability---but flags that
          the profile has no provenance (it was inserted by hand).
    \item \textbf{Parameter sensitivity:} The sign result
          $V_{\mathrm{core}}(0) \leq V_{\mathrm{core}}(q)$ is
          insensitive to $\sigma$, $\delta$, $L_0$, and $R$ for all
          elastic-model profiles.
\end{enumerate}

All numerical checks are tagged \epistemictag{Check}. They confirm the
analytic result but do not constitute independent evidence.

% ============================================================================
\section{Summary}
\label{sec:NMCP:summary}
% ============================================================================

\begin{resultbox}[title=NMCP WP2 Summary]
\textbf{Question asked:} Can any admissible physical mechanism within
$S_{\mathrm{EH}} + S_{\mathrm{NG}}$ generate non-monotone core energy,
producing a secondary minimum in $V(q)$?

\textbf{Answer: No.}

\textbf{Key results:}
\begin{center}
\begin{tabular}{lll}
\toprule
\textbf{Item} & \textbf{Result} & \textbf{Tag} \\
\midrule
General sign theorem & $V_{\mathrm{core}}(0) \leq V_{\mathrm{core}}(q)$
for all elastic cores & [Dc] \\
$\Zthree \to \Ztwo$ transition & No provenance in
$S_{\mathrm{EH}} + S_{\mathrm{NG}}$ & [Dc] \\
Non-separable coupling & Sign extends to general case & [Dc] \\
Core topology change & No framework in
$S_{\mathrm{EH}} + S_{\mathrm{NG}}$ & [Dc] \\
Provenance test & Failed: no mechanism generates non-monotonicity
& [Dc] \\
\bottomrule
\end{tabular}
\end{center}

\textbf{Classification:} Bounded no-go within
$S_{\mathrm{EH}} + S_{\mathrm{NG}}$ under elastic brane model.

\textbf{Impact:} The NMCP loophole from N7 is now closed within the
elastic brane model class. Combined with the N7 monotone no-go, the
\emph{entire} core-profile sector (monotone and non-monotone) is
bounded: no core-energy contribution within $S_{\mathrm{EH}} +
S_{\mathrm{NG}}$ can generate a secondary minimum in $V(q)$.

\textbf{Updated candidate landscape:}
\begin{center}
\begin{tabular}{llll}
\toprule
\textbf{Candidate} & \textbf{Before NMCP} & \textbf{After NMCP}
& \textbf{Status} \\
\midrule
N1 Israel thin-junction & Bounded no-go & Unchanged & Dead end \\
N7 monotone core & Bounded insufficiency & Unchanged & Dead end \\
N7 non-monotone core (NMCP) & [OPEN] & Bounded no-go & Dead end \\
N2 bulk backreaction & Bounded insufficiency & Unchanged & Dead end \\
\bottomrule
\end{tabular}
\end{center}

\textbf{Remaining escape routes (all non-minimal):}
\begin{itemize}
    \item Higher-order brane terms (DBI, Gauss--Bonnet)
    \item Topological contributions (Chern--Simons)
    \item Form-field coupling
    \item Violation of assumption~(D) by non-elastic core physics
\end{itemize}

\textbf{Anti-smuggling compliance:}
\begin{itemize}
    \item[$\checkmark$] No forbidden donors imported
    \item[$\checkmark$] No phenomenological profile inserted
    \item[$\checkmark$] No target-driven parameter selection
    \item[$\checkmark$] Provenance test applied before numerical work
    \item[$\checkmark$] All candidate mechanisms tested against WP1 criteria
    \item[$\checkmark$] Bounded no-go documented honestly
\end{itemize}
\end{resultbox}
